\chapter{Introduction}

\section{Introduction}

Social media platforms have fundamentally altered the landscape of human communication. Platforms like Twitter, Reddit, and Instagram have become spaces where individuals share thoughts, sentiments, and reactions in real time. Daily, millions of textual updates are generated across these networks, ranging from clear expressions of joy or sadness to complex, sarcastic, or inherently ambiguous discourse.

Analyzing the emotional content of social media discourse is valuable for diverse applications. Corporations seek to monitor consumer sentiment regarding products, while researchers investigate mental health trends through the analysis of online expressions. Consequently, there is a critical need for automated systems capable of processing vast textual corpora to identify underlying emotional states.

This project proposes a framework for the classification of social media text into four distinct emotional categories: Positive, Negative, Neutral, and Ambiguous. The inclusion of the latter category distinguishes this approach from conventional systems that lack explicit uncertainty handling.

\section{Background and Motivation}

Emotion detection from text is an established research problem that has been investigated for several decades. However, a significant majority of existing systems share a fundamental limitation: the requirement to produce a definitive classification for every input. Conventional models process inputs and map them to a single class within a fixed label set. This behavior persists even when the input text is genuinely unclear, forcing a potentially misleading classification rather than acknowledging uncertainty.

Consider the sentence ``I am fine.'' On the surface, it appears positive; however, in many social contexts, it is used sarcastically to convey the opposite meaning. A standard classifier trained on surface-level features will typically return \textsc{positive} for this sentence, failing to account for its frequent sarcastic usage. Such a prediction is inherently misleading, as it overlooks the contextual nuances of the expression.

This research is founded on the principle that uncertainty is a valid and informative output. A robust system should explicitly acknowledge instances of uncertainty, particularly in emotion detection where the cost of an erroneously confident prediction may exceed that of flagging the input as unclear.

A second primary motivation for this research is the utilization of geometric data structures to quantify ambiguity. When text is converted into a vector representation using a language model, semantically similar inputs are mapped to proximal regions in the embedding space. This project investigates the geometric structure of the neighbourhood around each input in the embedding space. If the neighbourhood is well-defined and homogenous, the input is likely easily discriminable. Conversely, if the neighbourhood exhibits complex geometric characteristics, the input likely resides in a boundary region between emotional clusters, indicating inherent ambiguity.

To measure this geometric complexity, the project uses Topological Data Analysis (TDA), specifically a technique called persistent homology. The application of persistent homology to detect ambiguity in social media sentiment represents a novel contribution to the field of affective computing. Recent studies on emotion topology~\cite{emotion_topology} have explored the fundamental components of emotions from text using word embeddings, but the application of local TDA for ambiguity detection remains a novel frontier.

A third motivation arises from an analysis of the performance of existing models on the
GoEmotions dataset. Models such as BERT-based classifiers and SVM systems achieve
high performance on majority classes while struggling with minority classes, and
typically lack the capacity to express predictive uncertainty. The GoEmotions dataset
is also known to exhibit significant label noise and ambiguous samples due to
disagreements among annotators. Furthermore, standard models often fail to capture nuances such as sarcasm~\cite{sarcasm}, leading to misleading results. Traditional models treat all inputs as equally discriminable, a limitation this project seeks to rectify through explicit ambiguity detection.

\section{Problem Statement}
 
The research problem addressed in this study is defined as follows. Given a short piece of text
from a social media platform, the objective is to build a system that:
 
\begin{enumerate}
    \item Classifies the text into one of three emotion categories: Positive, Negative, or
    Neutral.
    \item Determines whether the text is too ambiguous to be classified confidently, and if so,
    returns an Ambiguous label instead.
\end{enumerate}

The challenge is not merely building an accurate classifier, but determining when a
classification should not be made. Standard machine learning models are designed to optimize
accuracy over a fixed set of classes and therefore always produce a definitive prediction, even
when the input is inherently unclear.

Such models lack a mechanism to represent uncertainty or ambiguity in human language. This
project addresses that limitation by introducing a system that is capable of identifying and
explicitly handling ambiguous emotional expressions, rather than forcing potentially misleading
classifications.

\section{Objectives}

The primary objectives of this research are as follows:

\begin{enumerate}
    \item Build a sentiment classification system that works on informal social media text from
    the GoEmotions dataset.
    \item Use Topological Data Analysis, specifically Vietoris-Rips persistent homology, to
    measure the geometric complexity of the neighbourhood of each input in the embedding
    space.
    \item Combine topological complexity with classifier confidence to produce a principled
    ambiguity signal.
    \item Map the 28-class GoEmotions labels to three classes (Positive, Negative, Neutral)
    in a way that preserves the semantic groupings.
    \item Handle the class imbalance in the dataset without using synthetic data generation
    techniques that would distort the topological structure.
    \item Deploy the final system as an interactive command-line tool that explains its decisions.
\end{enumerate}

\section{Scope}

The system works with English text. It is trained on Reddit comments from the GoEmotions
dataset. It returns one of four outputs: Positive, Negative, Neutral, or Ambiguous. The system
is designed to run on a standard laptop CPU for inference, though training benefits from GPU
availability.

The system does not handle audio, images, or video. It does not track user history or learn from
new inputs after training. It operates as a single-query system where each input is classified
independently.

\section{Significance and Contribution}

The key contribution of this project is the use of \textbf{persistent homology as an ambiguity
signal at inference time}. When a new text arrives, the system finds its nearest neighbours in
the training embedding space and computes a fresh Vietoris-Rips persistence diagram from
those neighbours. The number of one-dimensional holes (H1 features) in that diagram measures
how geometrically tangled the neighbourhood is. A high count means the input is in a
complicated part of the space, which the system interprets as a sign of ambiguity.

This is different from existing approaches to uncertainty in NLP in two ways. First, it is
structural rather than just probabilistic. Most uncertainty methods rely on dropout-based
sampling or ensemble disagreement. This project uses the actual geometry of the data. Second,
it is computed fresh at inference time using real training data, not a precomputed approximation.

The approach is also interpretable. When the system flags a text as Ambiguous, it can explain
exactly which rule triggered the decision: extreme loop count, moderate loop count combined
with low confidence gap, or very low confidence gap alone. This kind of explanation is useful in
practice.

\section{Report Structure}

The remainder of this report is organized into seven chapters, each addressing a specific aspect of the project's development and evaluation:

\begin{itemize}
    \item \textbf{Chapter 2: Literature Survey} reviews the foundational research in sentiment analysis, emotion detection, handling uncertainty in NLP, and the application of Topological Data Analysis to text.
    \item \textbf{Chapter 3: Requirements Analysis} outlines the functional and non-functional requirements of the proposed system, including data specifications and design constraints.
    \item \textbf{Chapter 4: System Design and Architecture} details the high-level architecture, component interactions, and the data flow through the emotion detection pipeline.
    \item \textbf{Chapter 5: Implementation Details} provides an in-depth explanation of the technical implementation, including embedding strategies, topological feature extraction, and the project development timeline.
    \item \textbf{Chapter 6: Results and Evaluation} presents the experimental results, performance metrics on the GoEmotions dataset, and qualitative case studies.
    \item \textbf{Chapter 7: Conclusion and Future Work} summarizes the findings of the project, discusses the identified limitations, and outlines potential directions for future research.
\end{itemize}
